\documentclass[border=10pt,tikz]{standalone}
\usepackage{fontspec}
\setmainfont{Apple SD Gothic Neo}
\setsansfont{Apple SD Gothic Neo}
\usepackage{tikz}
\usetikzlibrary{positioning,arrows.meta,shapes,fit,backgrounds,calc,decorations.pathreplacing,decorations.pathmorphing}
\begin{document}
% _design.tex — shared TikZ design system for all blog diagrams
% Inputs nothing on its own — meant to be \input{} from each diagram .tex
% AFTER \begin{document}, BEFORE \begin{tikzpicture}.
%
% Provides a unified visual language across GoF / DSA / EMC++ / Beautiful series:
%   - Colors: muted, accessible, color-coded by role
%   - Node styles: consistent sizing, rounded, subtle borders
%   - Edge styles: UML-style inheritance, aggregation, plus dashed alternates
%
% Usage:
%   \documentclass[border=10pt,tikz]{standalone}
%   \usepackage{tikz}
%   \usetikzlibrary{positioning,arrows.meta,shapes,fit,backgrounds}
%   \begin{document}
%   \input{../_design.tex}     % adjust relative path
%   \begin{tikzpicture}[blog]  % apply 'blog' style
%   ...

\definecolor{absfill}{RGB}{220,232,247}   % cool blue
\definecolor{absbord}{RGB}{72,118,182}
\definecolor{confill}{RGB}{249,222,222}   % warm red
\definecolor{conbord}{RGB}{197,93,93}
\definecolor{impfill}{RGB}{223,238,222}   % soft green
\definecolor{impbord}{RGB}{99,154,93}
\definecolor{clifill}{RGB}{251,239,205}   % pale yellow
\definecolor{clibord}{RGB}{200,167,72}
\definecolor{hifill} {RGB}{253,224,200}   % accent orange
\definecolor{hibord} {RGB}{210,135,68}
\definecolor{nodefill}{RGB}{250,250,250}
\definecolor{nodebord}{RGB}{120,120,120}
\definecolor{edgec}  {RGB}{80,80,80}
\definecolor{edgesc} {RGB}{145,145,145}
\definecolor{groupbord}{RGB}{200,200,200}

\tikzset{
  blog/.style={
    font=\sffamily\small,
    every node/.style={font=\sffamily\small}
  },
  % --- node roles ---
  cls/.style={
    rectangle, rounded corners=2pt,
    draw=nodebord, line width=0.4pt,
    minimum height=8mm, minimum width=28mm,
    inner sep=3pt, fill=nodefill, align=center
  },
  abs/.style={cls, draw=absbord, fill=absfill, font=\sffamily\itshape\small},
  con/.style={cls, draw=conbord, fill=confill},
  imp/.style={cls, draw=impbord, fill=impfill},
  cli/.style={cls, draw=clibord, fill=clifill},
  hi/.style={cls, draw=hibord, fill=hifill, font=\sffamily\bfseries\small},
  % --- circle (for trees / graph nodes) ---
  vx/.style={
    circle, draw=absbord, line width=0.4pt,
    minimum size=8mm, fill=absfill, inner sep=0pt
  },
  vxc/.style={vx, draw=conbord, fill=confill},
  vxh/.style={vx, draw=hibord, fill=hifill, font=\sffamily\bfseries},
  % --- decision diamond ---
  q/.style={
    diamond, draw=clibord, line width=0.4pt, fill=clifill,
    inner sep=2pt, aspect=2.4, align=center
  },
  % --- edges ---
  inh/.style={                                  % UML inheritance
    -{Triangle[length=3mm,width=3mm,open]},
    draw=edgec, line width=0.5pt
  },
  agg/.style={                                  % UML aggregation
    {Diamond[length=2.5mm,width=2.5mm,open]}-{Stealth[length=2.2mm]},
    draw=edgec, line width=0.45pt
  },
  uses/.style={                                 % directed reference
    -{Stealth[length=2.2mm]}, draw=edgec, line width=0.5pt
  },
  arr/.style={                                  % plain arrow (graphs/trees)
    -{Stealth[length=2mm]}, draw=edgec, line width=0.45pt
  },
  edge/.style={                                 % undirected (trees)
    draw=edgec, line width=0.45pt
  },
  alt/.style={                                  % dashed: similar/alt
    -{Stealth[length=2mm]}, dashed, draw=edgesc, line width=0.45pt
  },
  ret/.style={                                  % dashed creates/returns
    -{Stealth[length=2mm]}, dashed, draw=edgesc, line width=0.45pt
  },
  thr/.style={                                  % threading link (red dashed)
    ->, dashed, draw=conbord, line width=0.5pt
  },
  % --- group / region ---
  group/.style={
    rectangle, rounded corners=4pt,
    draw=groupbord, line width=0.4pt,
    inner sep=10pt, fill=black!2
  },
  glabel/.style={font=\sffamily\bfseries\small, anchor=north west, inner sep=2pt}
}

% Korean font convenience macro — included by every Hangul diagram
\providecommand{\KOR}{}

\begin{tikzpicture}[blog]

% Axes
\draw[arr, line width=0.6pt] (0, 0) -- (7.5, 0) node[anchor=west, font=\sffamily\footnotesize] {Utilization (\%)};
\draw[arr, line width=0.6pt] (0, 0) -- (0, 5) node[anchor=south, font=\sffamily\footnotesize] {응답 시간 (R/T)};

% X-axis labels
\foreach \x/\val in {1.4/20, 2.8/40, 4.2/60, 5.6/80, 7/100} {
  \draw[line width=0.3pt] (\x, -0.1) -- (\x, 0.1);
  \node[font=\sffamily\scriptsize, anchor=north] at (\x, -0.15) {\val};
}

% Exponential curve (hockey stick)
\draw[line width=0.8pt, draw=conbord, smooth]
  plot coordinates {
    (0.7, 0.5)
    (1.4, 0.5)
    (2.1, 0.5)
    (2.8, 0.55)
    (3.5, 0.65)
    (4.2, 0.9)
    (4.9, 1.3)
    (5.6, 2.0)
    (6.3, 3.2)
    (7.0, 4.5)
  };

% Data points
\foreach \x/\y in {0.7/0.5, 1.4/0.5, 2.1/0.5, 2.8/0.55, 3.5/0.65, 4.2/0.9, 4.9/1.3, 5.6/2.0, 6.3/3.2, 7.0/4.5} {
  \fill[conbord] (\x, \y) circle (2pt);
}

% Example marker only; this is an illustrative curve, not a universal threshold
\draw[dashed, line width=0.5pt, draw=warnbord] (5.6, 0) -- (5.6, 4.8);
\node[font=\sffamily\footnotesize\bfseries, text=warnbord, anchor=south] at (5.6, 4.8) {예시 80\%};

% Illustrative latency-growth annotation
\node[font=\sffamily\footnotesize\bfseries, text=warnbord, anchor=east] at (5.45, 3.4) {지연 증가};

% Caption
\node[font=\sffamily\footnotesize, text=mutec, align=center, anchor=north] at (3.5, -0.8) {%
대기 시간은 시스템과 workload에 따라 달라집니다. 이 곡선과 80\% 표시는 예시입니다};

\end{tikzpicture}
\end{document}
